%-----------PROJECTS (USER)-----------
% Your real projects content lives here.
% Toggle rendering in config/resume-config.tex via \TemplateModetrue / \TemplateModefalse.

%-----------PROJECTS-----------
% Add entries using:
%   \resumeProjectHeading{\textbf{Name} $|$ \emph{Tech}}{Dates}
% Optional links: \ProjectLinksAuto{<github>}{<demo>} (use {} to omit).

\section{Projects}

\resumeSubHeadingListStart

% Optional links examples (use either, or both; remove any by replacing with {}):
%   \ProjectLinksAuto{https://github.com/USERNAME/REPO}{https://demo.example.com}
%   \ProjectLinks{https://github.com/USERNAME/REPO}{}
%   \ProjectLinks{}{https://demo.example.com}
\resumeProjectHeading
    {\textbf{Project Name}\ProjectLinksAuto{https://github.com/USERNAME/REPO}{https://demo.example.com} $|$ \emph{Tech 1, Tech 2, Tech 3}}
    {Mon YYYY -- Mon YYYY}

    \resumeItemListStart
        \resumeItem{One-line description of the project and what it does}
        \resumeItem{Hard thing you implemented (scaling, ML model, caching, auth, streaming, etc.)}
        \resumeItem{Impact (users, stars, latency improvement, revenue, accuracy, etc.)}
    \resumeItemListEnd

\resumeProjectHeading
    {\textbf{Project Name}\ProjectLinksAuto{}{} $|$ \emph{Tech 1, Tech 2}}
    {Mon YYYY -- Mon YYYY}

    \resumeItemListStart
        \resumeItem{What you built / contributed}
        \resumeItem{Technical detail / architecture}
        \resumeItem{Outcome / impact}
    \resumeItemListEnd

\resumeSubHeadingListEnd
